\paragraph{Finite existence certificates against $c_2=2$.}
Table~\ref{tab:two-orbit} gives exact half-rate certificates for
Lemma~\ref{tou:finite}. Every row has positive $Q_0-B_0$, is strictly
below Elias, and beats the prescription even with $c_1=1,c_2=2$.
The last row places the exceptional point $86/87$ of the gap outside
the tested radius. These are existence certificates for a suitable
parameter and domain; no particular domain is specified by the table.

\begin{table}[tb]
\centering
\begin{tabular}{rrrrrr}
\toprule
$p=2^b-1$: $b$ & $d$ & $n$ & $K$ & separation$/\eta$ & ratio $>$\\
\midrule
521 & 2 & 346 & 173 & $4/5$ & $2^{9}$\\
1279 & 3 & 862 & 431 & $6/7$ & $2^{13}$\\
9689 & 5 & 6638 & 3319 & $10/11$ & $2^{85}$\\
23209 & 7 & 16546 & 8273 & $14/15$ & $2^{119}$\\
44497 & 19 & 36060 & 18030 & $38/39$ & $2^{9}$\\
216091 & 43 & 190144 & 95072 & $86/87$ & $2^{6}$\\
\bottomrule
\end{tabular}
\caption{Unique-witness two-orbit existence certificates. The ratio is
$J/(n2^{2H_2(1/2)/\eta})$, with $\eta=(2d+1)/n$.
The table does not claim efficient sampling or decoding.}
\label{tab:two-orbit}
\end{table}

The certificates use $(m,D,c)=(171,87,0)$, $(285,144,1)$,
$(1325,664,3)$, $(2361,1182,5)$, $(1895,949,17)$, and
$(4419,2211,41)$, respectively.
They replace $\sqrt p$ by its integer ceiling and check all bounds
using integer arithmetic. Lucas--Lehmer establishes each prime;
$(2d+1)(b-1)>n$ gives a strict sufficient Elias certificate.
Using $J_0=\lceil\binom mD/\sqrt{D(m-D)(m+1)+1}\rceil$
from Lemma~\ref{tou:concentration}, a displayed exponent $w$ is certified by
\[
 J_0^{2d+1}>n^{2d+1}2^{2n+w(2d+1)}.
\]
An independent replay verifies these inequalities.
Small-field checks separately enumerate every parameter in two fixtures
over the field of order $1{,}000{,}003$, verify the excluded-parameter counts, and
classify every potentially nearby codeword on selected lines by
interpolation. Integer fixtures also check actual radical domains,
the conservative norm bound, and greedy witness recovery.
These smaller fixtures test the geometry and proofs; they are not
additional violations of the numerical prescription.

\paragraph{A length-$226$ half-rate certificate.}
An exact index-sum count gives a shorter example over the $338$-bit prime
\[
 p=204255\cdot2^{320}+1,\qquad n=226,\qquad K=113.
\]
Primality has the compact certificate
$11^{(p-1)/2}\equiv-1\pmod p$: every prime divisor of $p$ must then
be $1$ modulo $2^{320}$, and hence exceed $\sqrt p$.
For $m=111$, $D=57$, the index-sum class $S=3192$ contains exactly
\[
 J=441734689284929317660654122848
\]
supports. This is the coefficient of $y^{57}x^{3192}$ in
$\prod_{i=1}^{111}(1+yx^i)$, replayed by two independent recurrences.
The finite existence inequality is positive, so some domain over this
explicit field has these uniquely nearby parameters, gap $5/226$,
and far separation $4\eta/5$. Exact integer comparisons give
\[
 \frac{J}{n2^{2/\eta}}>\frac{119}{100},\qquad
 \frac{J}{n2^{1/\eta}}>2^{45}.
\]
The second inequality concerns the original $c_1=c_2=1$ prescription.
This certificate specifies the field and the support class, but not a
successful domain parameter $q$.

At rate $19/44$, a separate certificate has
$p=121\cdot2^{320}+1$, $n=220$, $K=95$, and $\eta=1/44$.
Here $(m,D,S)=(108,48,2616)$ and the exact class contains
$31495815243247545087807036124$ supports. The ratio exceeds $21/20$
for $c_2=2$ and $2^{43}$ for $c_2=1$, with $c_1=1$ in both cases.
Primality is certified by $3^{(p-1)/2}\equiv-1\pmod p$.
The non-half-rate comparison is also entirely integral, using
\[
 2^{nH_2(K/n)}=\frac{n^n}{K^K(n-K)^{n-K}}.
\]
The same existence and domain-specification qualification applies.
